% Rank-three denominator data for the Gritsenko--Nikulin Delta-five algebra.
% Primary anchors: Gritsenko--Nikulin 1998, Gritsenko 1995,
% Eichler--Zagier 1985, Borcherds 1998.

\providecommand{\CC}{\mathbb{C}}
\providecommand{\ZZ}{\mathbb{Z}}
\providecommand{\RR}{\mathbb{R}}
\providecommand{\NN}{\mathbb{N}}
\providecommand{\Delaf}{\Delta_{5}}
\providecommand{\LamII}{\Lambda^{2,1}_{\mathrm{II}}}
\providecommand{\gbkm}{\mathfrak{g}_{\Delta_{5}}}
\providecommand{\smult}{\operatorname{smult}}
\providecommand{\CoHA}{\operatorname{CoHA}}
\providecommand{\Pii}{\mathcal{P}_{\mathrm{II}}}
\providecommand{\Piprim}{\mathcal{P}_{\mathrm{II},\mathrm{prim}}}
\providecommand{\kBKM}{\kappa_{\mathrm{BKM}}}

\section*{Rank-Three Borcherds Denominator and \(K3\times E\) Separations}

\paragraph{Theorem 1: rank-three denominator.}
Let
\[
  \Lambda^{2,1}=\ZZ f_2\oplus\ZZ f_3\oplus\ZZ f_{-2},
  \qquad
  (f_2,f_{-2})=-1,\qquad (f_3,f_3)=2,
\]
with all other pairings zero. The Type-II real simple roots are
\[
  \delta_1=2f_2-f_3,\qquad
  \delta_2=2f_{-2}-f_3,\qquad
  \delta_3=f_3 .
\]
Their Gram matrix is
\[
  ((\delta_i,\delta_j))_{i,j=1}^{3}
  =
  \begin{pmatrix}
    2&-2&-2\\
    -2&2&-2\\
    -2&-2&2
  \end{pmatrix},
  \qquad \det=-32,
\]
with signature \((2,1)\). The chamber
\(\Piprim=\ZZ_{\geq0}\delta_1+\ZZ_{\geq0}\delta_2+\ZZ_{\geq0}\delta_3\)
is preserved up to permutation by the \(S_3\)-symmetry generated by the
reflections in \(f_2-f_3\) and \(f_{-2}-f_2\). Its three isotropic
vertices are
\[
  2f_2,\qquad 2f_{-2},\qquad 2(f_2-f_3+f_{-2}).
\]

The simple roots of the Borcherds--Kac--Moody superalgebra satisfy
\[
\begin{array}{ll}
(\alpha,\alpha)>0 &\Longleftrightarrow\ \alpha\in\Delta^{\mathrm{re}},\\
(\alpha,\alpha)\leq0 &\Longleftrightarrow\ \alpha\in\Delta^{\mathrm{im}},\\
\alpha\neq\alpha' &\Longrightarrow\ (\alpha,\alpha')\leq0,\\
\alpha\in\Delta^{\mathrm{re}} &\Longrightarrow\
  2(\alpha,\alpha')/(\alpha,\alpha)\in\ZZ .
\end{array}
\]
The final implication is the real-root Serre integrality condition.
Together these conditions promote the Lorentzian denominator data to a
Borcherds--Kac--Moody root system.

The Weyl vector is fixed by \((\rho,\delta_i)=-1\):
\[
  \rho=\frac12(\delta_1+\delta_2+\delta_3)
  =f_2-\frac12 f_3+f_{-2},
  \qquad
  (\rho,\rho)=-\frac32 .
\]
On the tube domain
\[
  \Omega(\LamII)=\LamII\otimes\RR+i\,\mathcal C(\LamII)_+,
\]
where \(\mathcal C(\LamII)_+\) is the chosen positive cone, the
Gritsenko--Nikulin denominator is
\[
  64^{-1}\Delaf(2Z)
  =
  e^{-2\pi i(\rho,Z)}
  \prod_{\alpha\in\Delta_+}
  \bigl(1-e^{-2\pi i(\alpha,Z)}\bigr)^{\smult(\alpha)}.
\]
The scalar \(64^{-1}\) fixes the monic product normalization; the weight
and root exponents are read from the Borcherds product.

Write the Eichler--Zagier generator as
\[
  \phi_{0,1}(\tau,z)=\sum_{a\geq0,\ell\in\ZZ} f(a,\ell)q^a r^\ell,
  \qquad
  f(a,\ell)=c(4a-\ell^2).
\]
For the root
\[
  \alpha(n,\ell,m)=2nf_2-\ell f_3+2mf_{-2},
  \qquad
  D=4nm-\ell^2,
\]
the denominator exponent is the signed supermultiplicity
\[
  \smult_{\gbkm}(\alpha(n,\ell,m))
  =
  \dim\gbkm_{\alpha,\bar0}-\dim\gbkm_{\alpha,\bar1}
  =
  f(nm,\ell)
  =
  c(D).
\]
Positive coefficients contribute even signed dimension; negative
coefficients contribute odd signed dimension. A parity-refined
root-space dimension requires the parity datum in addition to the
Jacobi coefficient.

The constant term of \(\phi_{0,1}\) is
\[
  \phi_{0,1}|_{q^0}=\zeta^{-1}+10+\zeta,
  \qquad c(0)=10.
\]
At a primitive isotropic vertex \(a_0\) of \(\Pii\), the even
imaginary-simple multiplicity is
\[
  \tau(a_0)=c(0)-c(-1)=10-1=9.
\]
The subtraction removes the real-root contribution. Equivalently,
the one-variable cusp identity reads
\[
  1-\sum_{t\geq1}m(ta_0)q^t
  =
  \prod_{t\geq1}(1-q^t)^{\tau(ta_0)},
\]
and the primitive exponent is the \(\eta^9\) factor in the
Gritsenko--Nikulin cusp expansion. The \(S_3\)-symmetry of the chamber
permutes the three isotropic vertices, so the value \(\tau(a_0)=9\)
is the same at each zero-dimensional cusp.

Borcherds' weight theorem gives
\[
  \kBKM(\Delaf)
  =
  \operatorname{wt}(\Delaf)
  =
  \frac{c(0)}2
  =
  5 .
\]
\emph{Proof.}
The Gram matrix follows by substitution from the three displayed
simple roots and the bilinear form on \(f_2,f_3,f_{-2}\). Solving
\((\rho,\delta_i)=-1\) gives the displayed Weyl vector. The product
formula is the Gritsenko--Nikulin Borcherds denominator attached to
the Eichler--Zagier input \(\phi_{0,1}\). The \(q^0\)-coefficient
\(\zeta^{-1}+10+\zeta\) gives \(c(0)=10\), and Borcherds' weight
formula gives \(\operatorname{wt}(\Delaf)=c(0)/2=5\). The same
coefficient gives \(\tau(a_0)=c(0)-c(-1)=9\) after subtracting the
real-root contribution \(c(-1)=1\).

\paragraph{Theorem 2: the four \(K3\times E\) scalars.}
Let \(Y=K3\times E\). The compact and specialised invariants are
\[
\begin{array}{c|c|l}
\text{invariant} & \text{value} & \text{construction}\\
\hline
\kappa_{\mathrm{cat}}(Y) & 0
  & \chi(\mathcal O_Y)=\chi(\mathcal O_{K3})\chi(\mathcal O_E)=2\cdot0\\
\kappa_{\mathrm{ch}}(Y) & 0
  & \sum_q(-1)^q h^{0,q}(Y)=1-1+1-1\\
\kappa_{\mathrm{ch}}^{\mathrm{Heis}}(Y) & 3
  & \kappa_{\mathrm{ch}}^{\mathrm{Heis}}(K3)+
    \kappa_{\mathrm{ch}}^{\mathrm{Heis}}(E)=2+1\\
\kBKM(\gbkm) & 5
  & c_1(0)/2=10/2\\
\kappa_{\mathrm{fiber}}(Y) & 24
  & \operatorname{rk}\widetilde H(K3,\ZZ)=1+22+1 .
\end{array}
\]
The set of values is \(\{0,3,5,24\}\), sourced by compact Hodge theory,
Heisenberg--Mukai specialisation, the Borcherds denominator, and the
Mukai lattice.

\emph{Proof.}
K\"unneth gives \(\chi(\mathcal O_{K3})=2\), \(\chi(\mathcal O_E)=0\),
and \(\chi(\mathcal O_Y)=0\). The Hodge diamond factorisation gives
\((h^{0,0},h^{0,1},h^{0,2},h^{0,3})=(1,1,1,1)\), hence
\(\kappa_{\mathrm{ch}}(Y)=0\). The Heisenberg--Mukai specialisation
has one rank from the elliptic factor and two ranks from the K3
holomorphic symplectic directions. The BKM value is Theorem 1. The
Mukai lattice rank is \(1+22+1=24\).

\paragraph{Theorem 3: primitive CHL denominator constants.}
For \(N\in\{1,2,3,4,6\}\) the primitive BKM denominator convention uses
the \(g_N\)-twined weight-zero index-one Jacobi input
\(\phi_{0,1}^{(N)}\) with discriminant-zero coefficient \(c_N(0)\):
\[
\begin{array}{c|ccccc}
N & 1 & 2 & 3 & 4 & 6\\
\hline
c_N(0) & 10 & 8 & 6 & 4 & 2\\
\kappa_{\mathrm{BKM}}(\Phi_N) & 5 & 4 & 3 & 2 & 1 .
\end{array}
\]
The Gritsenko--Clery diagonal-divisor rows have their own triples
\((t,N;k)\), cover groups, characters, and product normalisations. A
comparison with the primitive CHL row requires an explicit row map
recording all four data.

\emph{Proof.}
Borcherds' weight formula gives
\(\kappa_{\mathrm{BKM}}(\Phi_N)=\operatorname{wt}(\Phi_N)=c_N(0)/2\).
Gritsenko--Nikulin's CHL tower has weights \((5,4,3,2,1)\) at
\(N=(1,2,3,4,6)\), yielding the displayed constants. At \(N=1\),
\(\phi_{0,1}|_{q^0}=\zeta^{-1}+10+\zeta\), so \(c_1(0)=10\).

\paragraph{Square convention.}
The primitive \(N=1\) denominator consumes the Eichler--Zagier input
\[
  \phi_{0,1}(\tau,z)=\zeta^{-1}+10+\zeta+O(q)
\]
and produces the weight-five form \(\Delaf\). The K3 elliptic genus is
\[
  \operatorname{Ell}_{K3}(\tau,z)=2\,\phi_{0,1}(\tau,z).
\]
This doubled input belongs to the Igusa square convention:
\[
  \chi_{10}^{\mathrm{un}}=\Delaf^2,
  \qquad
  \Phi_{\!10}^{\mathrm{phys}}=(\Delta_{k(1)}^{(1)})^2 .
\]
Thus the primitive denominator weight is \(5\), while the square
convention has weight \(10\).

\paragraph{Operadic and Hall-algebra scope.}
For \(d\geq3\), the native CY-to-chiral output is
\[
  \Phi_d(\mathcal C)\in E_1\text{-}\mathrm{ChirAlg}.
\]
\(E_2\)-structures enter through centre, double, module-category, or
stable-envelope enhancements such as
\[
  Z^{\mathrm{der}}_{\mathrm{ch}}(A),\qquad
  D(Y^+(X)),\qquad
  \mathrm{Rep}^{E_2}(A).
\]
For the toric affine threefold,
\[
  \CoHA(\mathbb{C}^{3})=Y^+(\widehat{\mathfrak{gl}}_1).
\]
The vacuum algebra \(\mathcal W_{1+\infty}\) appears after adjoining
the opposite half and evaluating the centre. The \(K3\times E\)
Hall--Drinfeld branch is the automorphic Borcherds branch
\[
  \mathcal D_{\hbar}\bigl(\CoHA_{K3\times E}\bigr)\leadsto\gbkm .
\]
The Mukai self-mirror K3 Yangian branch has a separate current
presentation and a separate \(d=2\) origin.

\paragraph{Proof obligations for comparison theorems.}
A comparison with class \(\mathcal S\), a Hall product, a higher-level
CHL form, or a Yangian current presentation must exhibit:
\[
\begin{array}{ll}
\text{source} & \text{the chiral algebra and its operadic level};\\
\text{lattice} & \text{the Lorentzian lattice, simple roots, and Weyl vector};\\
\text{input} & \text{the Jacobi form, coefficient }c_N(0),\text{ and parity rule};\\
\text{product} & \text{the denominator identity and scalar normalisation};\\
\text{passage} & \text{the map from positive CoHA half to double, centre, or module layer}.
\end{array}
\]
These data are the minimum proof datum for an identification.
